\documentclass[10pt,a4paper]{amsart}
\usepackage[margin=19mm]{geometry}
\usepackage{amsmath,amssymb,mathtools}
\usepackage[scr=rsfs,cal=euler]{mathalfa}
\usepackage{xfrac}
\usepackage[unicode,hidelinks,bookmarks=false]{hyperref}
\newcommand{\ie}{i.e.\ }
\newcommand{\cf}{cf.\ }
\newcommand{\resp}{respectively\ }
\newcommand{\Subsection}{\subsection}
\providecommand{\nobreakdash}{\nobreak\hbox{-}}
\setlength{\emergencystretch}{2em}
\multlinegap0pt
\usepackage{bm}
\usepackage{bbm}
\usepackage{dsfont}
\usepackage{xspace}
\usepackage{cancel}
\usepackage{mathtools}
\usepackage{amsthm}
\makeatletter
\@ifundefined{thm}{\newtheorem{thm}{Thm}}{}
\@ifundefined{lem}{\newtheorem{lem}[thm]{Lemma}}{}
\makeatother
\makeatletter
\expandafter\ifx\csname SauveCompteurs\endcsname\relax
\newenvironment{SauveCompteurs}[1]{%
\newcommand{\monparametre}{#1}
\openexport{\monparametre_sauve}
  \Export{thm}\Export{section}\Export{subsection}\Export{subsubsection}
\closeexport}{}
\fi
\newcommand{\cp}{{\sf  p}}
\renewcommand{\label}[1]{} 
\expandafter\ifx\csname preuve\endcsname\relax
\newenvironment{preuve}[1][Preuve]{\begin{proof}[#1]}{\end{proof}}
\fi
\makeatother
\title[A Pansiot-type subword complexity theorem for automorphisms ]{A Pansiot-type subword complexity theorem for automorphisms of free groups}
\author{Arnaud Hilion, Gilbert Levitt}
\date{Source: arXiv:2208.00676v2 (CC BY 4.0)}
\begin{document}
\maketitle

\noindent\emph{Source and licence.} A Pansiot-type subword complexity theorem for automorphisms of free groups. Version arXiv:2208.00676v2 (CC BY 4.0), distributed under the Creative Commons Attribution 4.0 International licence, as recorded in the arXiv licence metadata for this exact version and shown on the abstract page https://arxiv.org/abs/2208.00676v2. Full rights notice: licenses/arxiv-2208.00676-CC-BY-4.0.txt.\par\smallskip

\noindent\textit{Editorial scope.} Counted unit: the Lem whose source label is \texttt{entour} in the article (source0.tex lines 978--986, source label \texttt{entour}), delivered together with the article's own complete original proof of it (source0.tex lines 988--1017, one \texttt{proof} environment of 30 lines). No mathematics is written, repaired or re-derived: statement and proof are the article's own text line for line. Required context retained uncounted: none. Every definition, display and label of those lines is retained unchanged. Omitted from this package and not redistributed: 1--977, 987--987, 1018--1736. Nothing inside a retained excerpt is omitted or altered: every retained body is delivered exactly as the article's lines are written, so no omission receipt of any kind accompanies this package. Citations: the retained excerpts carry no citation command at all, so no bibliography entry of the article is transcribed and no reference is redistributed. Reference adaptation: none was needed and none was made; every reference inside the delivered unit resolves inside it, so no unresolved reference or citation is produced. Printed slips kept literal: no printed word, symbol, hypothesis or number of the counted statement or of its proof was altered. Layout and class: the article's own class is \texttt{amsart} with packages including latexsym, amsmath, amsthm, amsfonts, amssymb, fontenc, inputenc, aeguill, url, enumerate, mdwlist, graphicx, xcolor, geometry; the wrapper is the lane's pinned \texttt{amsart} editorial wrapper at 10pt on A4, which restates the theorem-like environments the retained text needs and adds the article's own macro definitions that the retained text needs (\texttt{\textbackslash cp}, \texttt{\textbackslash label}), with extra wrapper packages (bm, bbm, dsfont, xspace, cancel, mathtools). The delivered numbering is the unit's own. Assets: no retained span contains a figure environment, a graphics-inclusion command, a PDF-page inclusion, a table or a TikZ picture; no image file is copied into this package, the package declares no asset dependency and no image file is redistributed with it. Licence: version arXiv:2208.00676v2 of this article is distributed under the Creative Commons Attribution 4.0 International licence, as recorded in the arXiv licence metadata for this exact version and shown on the abstract page https://arxiv.org/abs/2208.00676v2. A full rights notice is delivered at licenses/arxiv-2208.00676-CC-BY-4.0.txt. Characters: every retained excerpt is pure ASCII, so the delivered engine, XeLaTeX with its default Latin Modern text font, reports no missing character.
\begin{lem} \label{entour}
The total complexity $\cp_\delta(n)$ is at least quadratic 
if $\delta$ is an edge path with a complete splitting refining a splitting of one of the following forms:
\begin{itemize}
\item  
  $v\theta w$, where $\theta$ is an exceptional path and $v,w$ are growing;
  \item $v\theta w$, where $\theta$ is a linear edge with $f(\theta)=\theta u$,   the paths  $v,w$ are growing, and $w$ has at least two growing terms in its complete splitting.
\end{itemize}
\end{lem}

\begin{proof}
In the first case we write  $\theta=eu^p\bar e'$, with $f(e)=eu^d$ and $f(e')=e'u^{d'}$.
We then have  $f^k_\sharp(\delta)=v_keu^{a_k}\bar e'w_k$, where $ | eu^{a_k}\bar e' | $ is an affine function $ak+b$  with $a>0$ and   $ | v_k | , | w_k | $  grow at least like $ck$ for some $c>0$. Given $n$, consider, for all $k$ such that $n\ge ak+b\ge n/2$ and $ ak+b+cn \ge n$, 
the subpaths $\eta$ of length $n$ of $f^k_\sharp(\delta)$ containing the subpath $eu^{a_k}\bar e'$. Assuming, as we may, $c<\frac12$, and neglecting $b$, we see that $ak$ varies between $n-cn$ and $n$.

 We claim that the number of subpaths thus obtained is equivalent to $n^2$.    
 Indeed, since $ | v_k | , | w_k | $ grow at least like $ck$, there is enough space on either side of $eu^{a_k}\bar e'$  to construct, for each $k$, at least $n-ak$ subwords of $f^k_\sharp(\delta)$ containing $eu^{a_k}\bar e'$.  We thus get  a lower bound $\displaystyle\sum_{   {n-cn}\le ak\le{  n}}(n-ak) $, which is quadratic in $n$.




 Moreover,
%
the paths that we constructed are all distinct as abstract paths: $\eta$ determines $k$  because the subpath $u^{a_k}$ may be characterized as the only maximal subpath   of $\eta$ whose height   is less than $h(e)$ and $h(e')$, and whose length is at least $n/2$.   

 In the second case, we write $f^k_\sharp(\delta)=v_k\theta u^{ k} w_k$. We now consider $k$ such that $k | u | \ge n/2$ and subpaths $\eta$ of length $n$ containing $\theta u^k$. To show that they are distinct, we consider the maximal subpath   of $\eta$ whose height   is less than $h(\theta)$ and whose length is at least $n/2$. There is an added difficulty 
 because $u$ may be a prefix of $w_k$. 
 
 Write $w=xyz$ with $x$ a (possibly empty)  
   Nielsen path, $y$ a  growing term, and $z$ growing. Using induction on height, we may assume that $y$ is not a connecting path.  
    If it is an EG edge, or an exceptional path, or an   NEG edge with $f(y)=y\tau$, 
     every prefix of $w_k$ of the form $u^p$ is contained in $x$, 
 so distinctness holds. 
 

 The only case remaining is when $y$ is an NEG edge with $f(y)=\tau y$. Using induction, one may assume that  $\tau$ is a Nielsen path, so $f^k_\sharp(\delta)=v_k\theta u^{ k} x\tau^kyz_k$. Distinctness holds 
 unless $x\tau^k$ is a power of $u$ followed by a path of length less than $ | u | $. If this happens, we consider, for $k$ such that $k (| u |+| \tau | )\ge n/2$, the subpaths $\eta$ of length $n$ of $f^k_\sharp(\delta)$ containing $u^{ k} x\tau^k$. 
  As in the first case,  
 the number of these paths is equivalent to $n^2$ because $| v_k |$ and $|z_k|$   grow at least  linearly.
\end{proof}

\end{document}
